\section{Cost targets and the sector-storage bound}
\label{app:family-cost}
The new training, initial validation, confirmation, and reserved profiles use $h=1101$, $1111$, $1113$, and $1127$, respectively, in the vacuum, two-point, and scattering definitions of Appendices~\ref{app:coverage}, \ref{app:families}, and \ref{app:extended-coverage}. There are four one-mass vacuum loop orders, three fully massive vacuum orders, three massive two-point orders, massless two-point orders two through four, ladder orders one through three, and bubble-box orders two and three. The nonplanar box, planar double box, pentagon, and one-virtuality box complete the 22 labelled family-and-loop entries. All momenta and masses use the preceding definitions in a common squared-momentum unit.

Let $p$ be the number of physical propagators, $N$ the number of coordinates, and $\boldsymbol a^{(0)}=(1,\ldots,1,0,\ldots,0)$ the corner, with $p$ ones. For $j=0,1,2$, choose $i_j=1+((h+j)\bmod p)$. The dotted targets are $\boldsymbol a^{(0)}+2\boldsymbol e_{i_j}$. If $N>p$, choose $k_j=p+1+((h+j)\bmod(N-p))$ and use mixed targets $\boldsymbol a^{(0)}+\boldsymbol e_{i_j}-\boldsymbol e_{k_j}$. If there is no auxiliary coordinate, the mixed target has only the extra physical power. Remove duplicate vectors. These definitions explain the complete target list without assigning arbitrary linear combinations to a physical integral.

For batch $c$ and seed action $a$, let $T(c,a)$ be its measured complete time, including recovery. The regression target is
\begin{equation}
 y(c,a)=\log\frac{T(c,a)}{T(c,\mathrm{native})}.
 \label{eq:family-cost-label}
\end{equation}
For a fixed batch, the denominator is constant and positive, while the logarithm is strictly increasing. Therefore minimizing this target over $a$ gives the same minimum as minimizing $T(c,a)$. This equivalence motivates the label; finite training cannot guarantee the correct minimizer. Incomplete attempts receive a bounded-study penalty of at least 180 seconds, or three times their measured cost when larger. A model predicts $\widehat y(c,a)$, and the application changes the fixed action $a_0$ only if the best predicted action satisfies $\widehat y(c,a)<\widehat y(c,a_0)-\delta$. Validation selects $\delta\in\{0,0.03,0.08\}$. The model families use maximum depths two, four, and six; forest models have 32 trees and shared-model leaves contain at least two training examples. Where two complete training batches are available, a family-specific cost tree uses one split and one example per leaf, using only target-complexity features that do not depend on line labels. Its two measured training batches determine the two action-cost vectors. Separate initial validation and confirmation determine whether this small model transfers to new line positions. Live prediction uses the fitted tree parameters and does not import the training library.

The ten seed actions and truncation conventions are those of Appendix~\ref{app:learned-bounds}. Each distinct configuration is measured once per repetition; identical bounds and truncation share that measured record. Validation uses two repetitions, and reserved testing uses five. The historical training-only records use their original internal application timer; new labels use fresh-process time. Their common quantity is complete reduction cost with retries, but historical import cost is absent. Separate fresh-process validation and testing assess the deployed choice. Reserved timers exclude the study-wide administrative storage audit, which runs before each timed process. Training and validation records retain this common bookkeeping cost; all reserved controls use the same reduced recording overhead. Each reported interval resamples the available paired repetition indices within every matched request (five in both comparisons) and recomputes the ratio of summed medians. The 2.5 and 97.5 percentiles of 2,000 draws describe repeat variation for that workload.

The final test uses $h=1151$ and the same target construction, after the target-pattern gate has been fixed from confirmation costs alone. The gate compares maximum extra positive power and maximum numerator rank. It changes only deployment; it does not refit the cost tree.

For the storage change, write a sector as the bit mask $s=\sum_{i:a_i>0}2^{i-1}$. Let $\mathcal S$ be the declared top sectors. The already restricted initialization admits only sectors $s$ with $s\mathbin{\&}t=s$ for some $t\in\mathcal S$. Every set bit of $s$ is then a set bit of $t$, which gives
\begin{equation}
 0\le s\le t\le\max\mathcal S.
\end{equation}
Consequently an array with $1+\max\mathcal S$ entries contains every admitted sector index. Entries beyond this bound are never populated by the restricted initializer. Reducing allocation and the corresponding empty-entry scans preserves all admitted relations. This argument requires declared-sector initialization, no cut propagators, and no symbolic-IBP extension. The implementation retains native allocation otherwise. The tests use cold single-family reductions, reject symmetry records outside the allocated domain, and keep the complete propagator list for IBP generation and numerator transformations.

For the sector-order table, the tested domain has one top sector with the first $p$ bits set. Its admitted indices are exactly $0,\ldots,2^p-1$. We generate only these sectors, apply the unchanged propagator permutation and ordering keys, and sort them as before. The IBP derivative raises index $a_i$ only in a term proportional to $a_i$. A zero index gives no such term, and raising a negative integer by one cannot make it positive. Expressing scalar products in the denominator basis lowers other indices. Thus IBP generation from these seeds cannot introduce a positive auxiliary denominator. Removing unused sectors from a sorted list preserves the relative order of the retained sectors; the new rank labels therefore preserve comparisons between admitted integrals.

Raw equation templates can still contain terms that vanish after the seed indices are substituted. Before requesting an omitted sector weight, the implementation substitutes those indices and simplifies the corresponding coefficient exactly. It discards the term only if that coefficient is zero, and rejects any nonzero term outside the admitted domain. A bounded cache retains coefficient strings already proved to be zero, so repeated templates reuse that exact result. This check avoids assigning weights to absent terms while preserving every nonzero relation. The change reduces the ordering-table size from $2^N$ to $2^p$. General multi-top-sector inputs retain the original ordering table.

The minimum number of lines in a nontrivial sector also follows from a finite count. There are $\binom{N}{k}$ sectors with $k$ lines. If $z_k$ of them belong to the known trivial set, a nontrivial sector with $k$ lines exists exactly when $z_k<\binom{N}{k}$. The smallest such $k$ equals the minimum found by enumerating every sector. The recurrence $\binom{N}{0}=1$ and $\binom{N}{k+1}=\binom{N}{k}(N-k)/(k+1)$ follows by canceling factorials and supplies the counts without that enumeration. If every sector is trivial, the original initialization value is retained. These exact changes reduce storage and empty-sector work; they do not remove the native coordinate limit or the complexity of polynomial construction.

We separately rebuild the same implementation with Kira's supported 128-bit integral weights and repeat the storage comparison. Both sides of this separate comparison use that width. Wider integer labels preserve the ordering operations and do not change coefficient arithmetic. Their extra range addresses the integral-weight limit; it does not change the native limit of 31 integral coordinates.

The higher-loop storage profiles use $h=1131$, with $h=1132$ for four low-loop controls. The scalar full-coordinate test uses $h=1141$, the corner, a first-line dot, and one first-line plus two last-line dots. Its basis is the scalar corner, and all auxiliary coordinates remain zero. The independent coefficient recurrence is Eq.~\eqref{eq:extended-raising}.
